This chapter outlines the empirical strategies used across the four studies in this thesis. It demonstrates how established empirical techniques are applied to different retail settings to examine the effects of strategic decisions on customer behavior and performance outcomes. Studies 1–3 rely on transaction-level data from a grocery retailer loyalty program; Study 4 uses publicly available data from the Spotify music streaming platform. Across these contexts, quantitative modeling approaches are designed to quantify how strategic retail decisions influence customer behavior and resulting performance outcomes. Table~\ref{tab:empirical_overview} provides a summary of the four empirical studies, including context, data source, its coverage, method, and level of analysis for performance outcomes. Each section then describes the research setting, data and sample construction, modeling approach, and limitations related to both data and modeling choices.

\begin{table}[htbp]
\centering
\begin{tabular}{p{0.5cm} p{1.3cm} p{2cm} p{2.8cm} p{3.4cm} p{1.3cm}}
\hline
\textbf{Study} & \textbf{Context} & \textbf{Data source} & \textbf{Coverage} & \textbf{Method} & \textbf{Level of analysis} \\
\hline

1 & Multiformat retailing 
& Retailer scanner data from loyalty program 
& 135,776 shopping trips over five selected product categories based on 4,921 customers 
& Error-correction model of sales response with price elasticity estimation and subsequent profit elasticity 
& Weekly brand sales \\
\hline

2 & Multiformat retailing 
& Retailer scanner data from loyalty program 
& 897,743 shopping trips from 5,468 customers 
& Two-stage approach: Price elasticity estimation followed by moderator analysis 
& Weekly brand sales \\
\hline

3 & Multichannel retailing 
& Retailer scanner data from loyalty program 
& 116,982 shopping trips across 578 households (147 of which adopt an online channel) 
& Propensity score matching combined with difference-in-differences estimation 
& Weekly customer purchases \\
\hline

4 & Digital platforms 
& Spotify playlist API and web-scraped data 
& 39,918 playlists with 49 consecutive weeks of follower data 
& Playlist-level unit fixed effects panel regression 
& Weekly playlist followers \\

\hline
\end{tabular}
\caption{Overview of empirical contexts, data, and methods across studies}
\label{tab:empirical_overview}
\end{table}

\section{Multiformat Retailing (Studies 1 and 2)}
The first two empirical studies are situated within the context of multiformat grocery retailing. They examine how pricing strategies perform, across different store formats, for one multi-format retailer. Despite belonging to the same retail chain, different formats cater to distinct shopping goals and create a unique context to explore customer responses to pricing decisions across different physical stores.

The data comes from the loyalty program of a major grocery retailer that operates in three store formats: hypermarkets, supermarkets, and convenience stores. From large-format stores catering to stock-up trips to medium-sized stores serving regular household needs and small outlets focused on immediacy and proximity, these formats differ substantially in size, assortment, and price positioning. Because shopping missions, basket sizes, time constraints, and assortment breadth systematically vary across these formats, customers are likely to respond differently to regular prices and promotions. For example, price discounts may be more salient for large, planned baskets at hypermarkets than for quick, convenience-driven trips at small stores. Promotional campaigns are coordinated centrally, while individual store managers retain flexibility in setting regular prices based on local market conditions. Taken as a whole, the situation reflects a multi-format strategy that combines service-oriented and price-competitive elements. Rather than following a strict everyday-low-price policy, the retailer maintains competitive regular prices but also relies on centrally coordinated promotions. This strategy is consistent with a hybrid or moderated hi–lo pricing format that lies between pure everyday-low-price and pure hi–lo in grocery retailing. The dataset includes detailed transaction-level records for approximately 5,500 customers over 153 consecutive weeks (2014–2016), covering around 900,000 shopping trips. Each record contains product-level information, including brand, category, regular price, promotional discount, and quantity purchased. This level of granularity allows for detailed measurement of customer behavior and pricing outcomes. From this dataset, study-specific samples are selected based on the requirements of each empirical design, resulting in differences in sample size and composition between Studies 1 and 2 (see Table~\ref{tab:empirical_overview}).

To evaluate how pricing strategies affect sales and profitability across physical store formats, the first two studies model customer responses using aggregated brand-level sales data. Aggregate sales response models are commonly used in retail pricing research because pricing decisions are typically implemented at the brand or category level rather than at the individual-customer level. Using brand-week panels therefore aligns the empirical model with the level at which retailers observe performance and make pricing decisions. Compared with individual-level choice models, this approach also reduces computational complexity while still capturing meaningful variation in customer responses to prices and promotions across store formats (\cite{Pauwels2007-ky,Datta2022-fz}). Thus, the unit of analysis of these studies is at the brand-week level, aggregated from transaction data, allowing the studies to quantify market-level customer response without relying on individual-level behavioral tracking. This approach is especially suitable for retailer-facing questions, such as whether specific pricing strategies improve sales or profitability across different retail formats.

To align with this modeling approach, the two studies leverage the same dataset to construct comparable brand-level sales panels across the three formats. The analysis focuses on the three store formats belonging to the same retail chain, located in the same area, for the purpose of ensuring the consistency and comparability of formats. To avoid biases due to assortment variation, products were selected based on consistent availability (a) across all three formats and (b) throughout the observation period, with an emphasis on frequently purchased brands. To ensure that any observed pricing effects were not confounded by size-based price differences, only items with standardized package sizes were retained. This approach mitigates the risk that format-based variation in sales is driven by differences in what is stocked, rather than by how customers respond to pricing.

In the first study, the focus is on profit-oriented evaluation of regular price setting for multi-tier PLs. The modeling follows a multi-step regression framework. First, price elasticities are estimated for each PL tier and format using observed variation in prices paid over time. These elasticity estimates are then combined with cost and price data to calculate contribution margins and profit effects. This two-part approach enables direct assessment of which tier–format combinations are most profitable, offering actionable insights for PL portfolio management (e.g., \cite{Geyskens2010-bx,Keller2020-so,ter-Braak2013-uq}).

In the second study, the focus shifts to how reference prices and discount framing affect sales across store formats. The study uses a two-stage regression design: the first stage estimates brand-level sales responsiveness to regular prices and promotions, and the second stage tests how these elasticities are moderated by brand characteristics, category type, and store format. This allows for identification of conditions under which either regular price changes or promotional discounts are more effective with respect to brand and category factors (e.g., \cite{Datta2022-fz,Haans2011-cv,Pauwels2007-ky}).

Both studies apply a sales response modeling framework (e.g., \cite{Pauwels2007-ky}), capturing how changes in price influence weekly sales volumes. An error correction specification is used because it allows the model to separate immediate demand reactions from longer-run equilibrium adjustments, which is particularly important when prices and promotions vary frequently across weeks. To mitigate endogeneity concerns, both models incorporate a set of control variables and fixed effects. These include brand fixed effects, time dummies, and category-level controls that help account for unobserved factors such as brand popularity, demand seasonality, and promotional cycles (\cite{Datta2022-fz}).

This modeling strategy offers several advantages. First, it mirrors managerial practice, where performance metrics like sales and margins are typically monitored at the weekly brand or category level (\cite{Widdecke2022-xo}). Second, by aggregating across customers, the models avoid the need for individual-level tracking (which is often unavailable in practice) while still revealing meaningful market-level behavioral responses. Third, the multi-step and multi-stage design enables a clear separation between estimating price elasticity and translating those effects into strategic insights about pricing, promotions, and profitability across brands, categories and formats.

While the dataset focuses on loyalty card holders from a single retail chain, it offers a high degree of consistency and detail, allowing precise observation of behavioral and sales dynamics over time within a stable customer base. Several limitations should however be acknowledged. First, the data do not capture purchases made at competing retailers; this may limit the generalizability of findings beyond the focal chain. Second, the analysis focuses on products consistently available across formats and time, enhancing comparability but possibly excluding niche or newly introduced items that could behave differently. Third, as the data reflect pricing and assortment decisions within one retailer’s organizational structure, some format-specific effects may be influenced by chain-level strategies that are not directly observable. In addition, the aggregate structure cannot capture within-brand heterogeneity in customer responses—for example, how different shopper segments may react to a given promotion (\cite{Ailawadi2004-ig, Geyskens2010-bx}). Furthermore, observed price variation is not randomly assigned and may reflect broader retailer strategy, so caution is required in interpreting causal claims (\cite{Bijmolt2005-kg, Datta2022-fz}). Despite these limitations, the combination of careful sample construction, consistent product selection, and inclusion of fixed effects for products, time, and store formats helps mitigate potential confounds from unobserved heterogeneity. This provides a reliable basis for assessing how pricing strategies influence performance across store formats.

\section{Multichannel Retailing (Study 3)}
The third empirical study, focusing on how customer purchase behavior evolves following the introduction of an online shopping channel, is set in the context of multichannel grocery retailing. Like the previous studies, it relies on transaction data from the retailer’s loyalty program, but it shifts the focus from store-format pricing to customer-level behavioral changes following online channel adoption. This setting provides a quasi-experimental setting to examine long-term shifts in behavior among customers who adopt (i.e., start using) the new channel, particularly whether online channel adoption leads to incremental revenue or simply redistributes spending across formats. \footnote{According to \textcite{Goldfarb2022-jt},
Quasi-experiment refers to the use of an experimental mode of analysis and interpretation to data sets where the data-generating process is not itself intentionally experimental (Campbell 1965). Instead, quasi-experimental research uses variation that occurs without experimental intervention but is nonetheless exogenous to the particular research setting.
}

The data come from the loyalty program of the same grocery retailer assessed in Studies 1 and 2, covering a four-year period that spans both pre- and post-introduction periods of the online channel. The online channel was introduced in late 2015, and this retailer was the only grocery chain in the region offering such an option at the time. This exclusivity minimizes competitive interference and enhances the internal validity of observed behavioral changes. The dataset combines two consecutive panels (October 2014–November 2016 and August 2017–July 2019) and tracks 578 households (147 of which adopt the online channel) across 116,982 shopping trips and roughly 1.75 million purchased products. The dataset includes customer-level purchase histories across both online and offline purchases, allowing for precise tracking of when each customer began using the online channel and how their purchasing patterns changed over time. Each transaction contains information on date, product category, quantity, spending, and purchase channel.

To identify the causal effect of online channel adoption on purchasing behavior, the study employs customer-level panel modeling techniques based on propensity score matching (PSM) and difference-in-differences (DiD) estimation. This approach is appropriate because online channel adoption is a voluntary decision rather than a randomized treatment. Propensity score matching, to reduce observable selection bias, first constructs comparable groups of adopters and non-adopters based on pre-adoption shopping behavior. Difference-in-differences estimation then exploits the panel structure of the data to compare behavioral changes before and after adoption at the customer level, thus helping isolate the causal impact of channel adoption from common time trends (\cite{Boehm2008-it, Campo2015-pt}).

To align with this modeling approach, the study leverages the loyalty data’s panel structure to track individual customers before and after the online channel introduction. Customers who adopted the online channel are compared against those who never adopted it throughout the observation window. This structure allows for the separation of pre-existing differences from actual behavioral shifts following adoption, thereby mitigating potential endogeneity concerns. By observing a panel of customers across a multi-year window (covering the before and after of the online channel’s introduction), the study captures longitudinal changes in behavior that follow the introduction and adoption of this new shopping format (\cite{Avery2012-ex,Bilgicer2015-ng}). This quasi-experimental setting allows the model to distinguish between behavioral changes caused by adoption versus those driven by general time trends or unobserved customer characteristics (\cite{Li2015-tq}).

This modeling strategy offers several advantages. First, it allows for customer-level inference, capturing the heterogeneity in how individuals respond to additional format adoption (\cite{Kushwaha2013-iu}). Second, the panel structure enables robust within-customer comparisons, reducing bias from unobserved time-invariant traits such as household size or income (\cite{Konus2008-mp}). Third, the combination of PSM and DiD mitigate self-selection bias by ensuring that adopters and non-adopters are comparable based on observable characteristics before the channel was introduced (\cite{Li2015-tq}), while DiD specifically controls for time-invariant unobservables and shared temporal trends.

While the dataset offers a unique opportunity to study an effect of online channel adoption on customer purchase behaviors, several limitations should be noted. First, although the retailer was the sole provider of online grocery shopping at the time of online channel introduction, competing retailers began offering similar services toward the end of the observation period; these are not explicitly accounted for. Second, because purchasing activity varies across customers, the ability to model detailed temporal dynamics (e.g., behavioral changes after one, two, or three years of adoption) is limited. Third, as the data reflect the early phase of adoption, the number of online users is relatively small and likely skewed toward early adopters, which may affect generalizability. In addition, DiD relies on the parallel trends assumption, that in the absence of treatment, adopters and non-adopters would have followed similar trajectories over time (\cite{Li2023-ho}). This assumption cannot be directly tested and may become less plausible as the post-adoption period extends. While matching helps improve group comparability, the use of proxies such as diaper purchases and store distance may not fully capture deeper household characteristics. Furthermore, DiD is less effective when treatment timing varies across individuals, or when adoption reflects structural changes in customer behavior that evolve gradually rather than in discrete shifts (\cite{Roth2023-xe}). Despite these limitations, matching and DiD allow comparison of purchasing behavior between adopters and non-adopters over time. The combined use provides a transparent and meaningfully practical basis for assessing how online channel adoption affects customer behavior over time.

\section{Digital Platforms (Study 4)}

The fourth empirical study is situated in the context of digital platforms, where customer behavior can be largely shaped by platform structures and mechanisms such as algorithmic curation and editorial decisions. Because these systems determine what users see and engage with, they influence which participants on the platform gain attention, or traffic. In such environments, potential external disruptions like the COVID-19 pandemic can alter user behavior at scale, raising questions about how users interact with content shift after the onset of such an event. Unlike the other studies, which rely on retailer loyalty data, this one uses publicly available platform data to examine observable user responses in a digital marketplace.

The data are extracted from publicly available sources related to Spotify, one of the world’s leading music streaming platforms. The dataset combines two primary sources. First, playlist-level follower data was obtained from Chartmetric, a commercial analytics provider that aggregates historical data from Spotify’s web API. This dataset includes weekly follower counts, curator type (e.g., Spotify-owned, label-affiliated, or independent), and content features such as the share of tracks from major labels and average track popularity. Second, data from Every Noise at Once is used to identify weekly instances of playlists featured on that platform, which is equivalent to platform-driven promotion that increases visibility. By integrating this information, the analysis distinguishes between follower growth (arising from organic user responses) and platform-driven exposure. The starting dataset from Chartmetric included over 1.2 million playlists. From this, playlists were filtered in stages: first, only those with sufficient data continuity were retained, and those lacking key variables, such as curator type, genre, or content attributes, were excluded. The final sample consists of 39,918 active playlists with 49 consecutive weeks of data spanning October 2019 to October 2020, thus covering periods before and after the pandemic was declared.

To examine changes in user responses on the platform, the study models playlist follower growth using fixed effects panel regression. Weekly follower changes serve as a proxy for aggregate observable user responses. A fixed effects panel specification is used because playlists differ systematically in their baseline popularity, genre, and curator characteristics. By controlling for these time-invariant playlist attributes, the model isolates within-playlist changes in follower growth over time (\cite{Papies2022-ln}).

To align with this modeling approach, the dataset is structured as a playlist-level panel observed over time, thus allowing each playlist to be tracked before and after the pandemic declaration. Each playlist serves as an observation unit, with time-varying attributes such as follower counts, content composition, curator type, and promotional exposure. The resulting dataset enables comparison of changes in follower growth within the same playlist over time, rather than across inherently different playlists. Key independent variables include playlist curator type, content composition, and weekly indicators of platform-driven promotion (i.e., Search Page featuring).

This modeling strategy offers several advantages. First, the playlist-level panel allows for detailed tracking of engagement changes across a broad sample of Spotify playlists. Second, the fixed effects specification accounts for unobserved heterogeneity at the playlist level, such as theme, brand strength, or historical growth patterns, which supports more credible within-playlist comparisons over time. Third, the inclusion of week fixed effects helps control for seasonality and platform-wide shifts in user activity, allowing changes in follower growth to be more closely associated with the pandemic period rather than general time trends.

First among limitations to be noted with regard to the fourth study, the dataset captures publicly observable playlist-level metrics such as follower counts rather than individual user behavior, and this leads to limited insights concerning the micro-level drivers of engagement. Second, the influence of Spotify’s algorithmic or editorial decisions is not directly observable and must be inferred from the featuring data. Third, follower growth serves as a proxy for user response and does not capture listening intensity or financial outcomes. In addition, the fixed effects model might be better suited for estimating large, persistent changes in engagement rather than subtle or short-lived shifts. As such, for smaller effects (particularly those limited to a subset of playlists or time periods) it may be harder to isolate the effect with certainty and subject to attenuation. Finally, even with week fixed effects, any unobserved time-varying shocks that affect playlists differently may not be fully captured by the model structure. Despite these constraints, the panel structure and fixed effects specification allow for consistent measurement of changes in user responses over time, particularly in response to large-scale external disruptions such as the pandemic (\cite{Denk2022-rt,Sim2022-ns}).

